\section{Independent proof audit of the AX11--Q--6 composition}
\label{sec:ax11q6-audit}

This section proves the new command-level composition.  It separates the
inherited numerical-calculus calibration, the new proof compiler, the
enumerator, and the literal physical machine.  The ordinary mathematical
argument is formalizable in PA; this section is not a claim that a new
end-to-end proof-assistant theorem has been compiled.

\subsection{Numerical formulas and inherited calibration}

Put $p(a,b)=(a+b)(a+b+1)/2+b$ and, for a fixed $k\geq4$, put
\[
 I_k(a,b)=kp(a,b),\quad M_k(a,b)=kp(a,b)+1,\quad
 U_k(a,B)=kp(B,a)+2,\quad E_k(a,b)=kp(a,b)+3.
\]
The distinguished contradiction is $1=M_k(0,0)$, truth is $0$, and
$N_k(B)=I_k(B,1)$.  In particular $I_k(0,0)=0$; neither numerical formulas
nor numerical stacks are assumed to be free term algebras at zero.

The separated-diagonal calculus retains six logical heads, the restricted
implicational head, the diagonal head $\Gamma_k(B)=I_k(B,U_k(B,B))$, and
the guarded membership-vacuity/pointed-separation gate.  The Q version
replaces its sixth head
\[
 U_k(D,A)\to\bigl(U_k(D,A\to H)\to U_k(D,H)\bigr)
\]
by
\[
 U_k(D,A\to H)\to\bigl(U_k(D,A)\to U_k(D,H)\bigr).
\]
The displayed arrows are $I_k$.  The exchange template
\[
 (P\to(Q\to R))\to(Q\to(P\to R))
\]
uses only the positive implicational $K,S$ basis.  In the supplied
calibration that basis follows from the retained restricted implicational
head alone, before the sixth head is used.  Applying the exchange template
twice proves that the ordinary and Q versions have uniform
primitive-recursive derivation translations, preserving conclusions.
There is no circular use of the distribution axiom being replaced.

The inherited radix calibration establishes the PA-refutation equivalence
for each fixed $k\geq4$.  Its substantive ingredients must be retained:
an injective hereditary recoding of both variable labels and formulas;
the generic derivation of unrestricted implication instances before
recoding a restricted leaf; and a total PA interpretation of all numerical
codes, including the unused residues $4,\ldots,k-1$.  Restricting soundness
only to the image of the old radix-four syntax would be insufficient for
the actual exhaustive numerical search.  The restricted implication
order is not preserved by hereditary recoding and cannot be silently
transported as a numerical inequality.

More explicitly, set $F_k(0)=0$ and, for $n>0$ with
$n=4p(a,b)+r$, $0\leq r<4$, set
\[
                F_k(n)=kp(F_k(a),F_k(b))+r.
\]
Both $a,b$ are smaller than $n$, so this is a primitive-recursive
course-of-values definition.  A partial inverse is defined at
$m=kp(a,b)+r>0$ by
$J_k(m)=4p(J_k(a),J_k(b))+r$ for $r<4$, and $J_k(m)=0$ otherwise,
with $J_k(0)=0$.  Induction proves $J_k(F_k(n))=n$, hence injectivity.
Recoding the numerical variable labels by this same $F_k$ gives
\[
 F_k(U_4(a,B))=U_k(F_k(a),F_k(B)),\qquad
 F_k(\Gamma_4(B))=\Gamma_k(F_k(B)).
\]
Thus binding and shadowing are preserved, as are the distinctness guards
$c\ne a,b$ of the gate.  The map is an embedding, not a bijection of the
whole target numerical syntax.  Its failure of order preservation is real:
$40<44$ but $F_5(40)=75>55=F_5(44)$.  Each restricted implicational leaf
is replaced by the fixed \texttt{rho8All} template at its recoded parameters;
that template's own side conditions follow from
$a,b\leq p(a,b)\leq I_k(a,b)$.

For the converse syntactic translation, decode every positive target code
$n=kp(a,b)+r$.  Translate residues zero, one, two and three respectively
as $T_k(a)\to T_k(b)$, $\operatorname{Bit}(x_a,x_b)$,
$\forall x_b T_k(a)$, and $x_a=x_b$.  Translate zero and unused residues
as closed truth.  Every source label in $T_k(B)$ is below $B$, so the
diagonal head is a valid vacuity instance for arbitrary target numbers,
not merely embedded ones.  The exceptional equation $I_k(0,0)=0$ is
handled by a fixed proof of truth equivalent to truth implying truth.
For the guarded separation head, the inherited PA finite-comprehension
proof supplier uses the matrix
$(\operatorname{Bit}(x_a,z)\lor z=x_b)\land T_k(c)[z/x_b]$
when $a\ne b$.  When $a=b$, its matrix is instead
$(z=x_b)\land T_k(c)[z/x_b]$, since the old membership antecedent becomes
$\operatorname{Bit}(z,z)$ under that binder.  This distinction and the
freshness of $c$ prevent an invalid free-variable treatment of colliding
labels.  These are syntactic PA proof suppliers, not a uniform arithmetic
truth predicate or a reflection principle.

With those inherited results and the exchange template, if
$\operatorname{Deriv}_{k,Q}(d,1)$ is the resulting numerical derivation
predicate, the calibration boundary is
\[
 \mathrm{PA}\vdash
 \exists d\,\operatorname{Deriv}_{k,Q}(d,1)
 \quad\longleftrightarrow\quad
 \exists p\,\operatorname{Prf}_{\mathrm{PA}}(p,\ulcorner0=1\urcorner).
 \tag{Cal}
\]
The case used by the new machine is $k=6$.  This audit neither replaces
(Cal) by a standard-model consistency argument nor promotes the repository's
unfinished iPRA endpoints to completed theorems.

The exact retained dependencies for (Cal) are
\file{candidate/proof/references/radix_calibration.tex} and
\file{candidate/proof/references/q_calibration.tex}, together with their
retained original, diagonal, and separate-head calibration sources.
\file{candidate/PROVENANCE.json} binds each source to its pinned repository
path and SHA-256; the finite premise-exchange derivation is replayed by
\file{candidate/proof/exchange_certificate.py}.

\subsection{Exact command interface}

Write $[a;w]=p(a,w)$; its total projections satisfy
$\operatorname{hd}(0)=\operatorname{tl}(0)=0$.  The store is $(t,w,d,h)$:
visible top, encoded theorem stack, offset, and cached theorem.  The new AX
command has \emph{four} operands $(c,x,y,z)$.  Set $D=d+x$ and $H=h+y$.
Let $q_0$ be the guarded separation gate, or zero when its numerical guard
fails; let $q_1=\Gamma_k(c)$; let $q_2,\ldots,q_7$ be the six logical
heads in reverse stack order, with $q_2$ the Q head; let $q_8$ be the
restricted implicational head; and let $q_9=U_k(d,h)$.  Then
\[
 \mathsf A_{c,x,y,z}(t,w,d,h)
 =\bigl(0,[q_0,q_1,\ldots,q_9,t;w],z,0\bigr).       \tag{AX}
\]
There are nine axiom entries, the generalized cached formula $U(d,h)$,
and the saved old top, so eleven
entries are pushed, subject to the numerical zero alias.

The extra operand is consumed by falling through to the complete MP field
reader.  At that point the visible top, offset and cache are all zero.
For every $z$, $I_k(0,z)=0$ if and only if $z=0$.  The subsequent no-push MP
therefore returns visible zero and cached zero, with offset $z$, whether
its comparison succeeds or fails.  Thus (AX) holds for every fourth
operand, not just the zero operands chosen by the proof compiler.

The other commands, with $b=d+x$, are exactly
\[
 \mathsf S(t,w,d,h)=
 \begin{cases}
  \mathsf{halt},&t+h=1,\\
  (h,\operatorname{tl}(w),d,\operatorname{hd}(w)),&t+h\ne1,
 \end{cases}
\]
\[
 \mathsf{MP}_x(t,w,d,h)=
 \begin{cases}
  (b,w,b,0),&I_k(t,b)=h,\\
  (0,w,b,0),&I_k(t,b)\ne h.
 \end{cases}                                                \tag{MP}
\]
In particular the sum-based Shift guard is retained.  It does not halt at
$(t,h)=(1,1)$, and the proof must not treat it as a cached-head-only guard.

\subsection{Uniform stack operations, including contradiction inputs}

Every $q_j$ for $0\leq j\leq8$ is divisible by $k$, including a failed
guard's zero.  The generalization head $q_9$ is congruent to two modulo
$k$ and is at least two.  Directly applying the Shift equation therefore
gives, without an intermediate halt,
\begin{align*}
 \mathsf A_{c,x,y,0};\mathsf S^{10}(t,w,d,h)
   &=(q_8,[t;w],0,q_9),\\
 \mathsf A_{c,x,y,0};\mathsf S^{11}(t,w,d,h)
   &=(q_9,w,0,t).
\end{align*}
The first sums tested are sums of multiples of $k$; the final sums contain
$q_9\geq2$.  This explains why these identities also hold at arbitrary
$t$, including $t=1$.

With $A_0=\mathsf A_{0,0,0,0}$, define
\[
 \begin{aligned}
  \mathsf{Push}&=A_0;\mathsf S^{10};\mathsf{MP}_0,&
  \mathsf{Pop}&=A_0;\mathsf S^{12},\\
  \mathsf{Discard}&=\mathsf{MP}_0,&
  \mathsf{Offset}_i&=\mathsf{MP}_i.
 \end{aligned}
\]
For all displayed natural values,
\begin{align*}
 \mathsf{Push}(t,w,d,h)&=(0,[t;w],0,0),\\
 \mathsf{Pop}(t,[a;w],d,h)&=(t,w,0,a),\\
 \mathsf{Discard}(t,w,0,h)&=(0,w,0,0),\\
 \mathsf{Offset}_i(0,w,0,0)&=(0,w,i,0).
\end{align*}
For Push the final MP comparison fails by the distinct implication and
universal residues.  For Discard both branches yield zero.  For Offset the
zero case succeeds with conclusion zero; each positive case fails.
Pop's last Shift tests $q_9+t\geq2$, so Pop preserves a visible
contradiction instead of halting.  The derivation includes $[0;0]=0$ and
the empty stack; no assertion of a free-list representation is needed.

For $0\leq j\leq8$, the selector
\[
 \mathsf{Discard};\mathsf A_{A,D,H,0};
       \mathsf S^{j+2};\mathsf{Pop}^{9-j}                 \tag{Sel}
\]
sends $(t,w,0,0)$ to $(q_j,w,0,0)$.  The exact number of Pops accounts for
the generated suffix and the saved zero.  The restricted implicational
leaf with old offset $u\leq D$ is selected by
\[
 \mathsf{Discard};\mathsf{Offset}_u;
 \mathsf A_{A,D-u,H,0};\mathsf S^{10};\mathsf{Pop}.
\]
These are different selectors from those of a ten-entry AX interface.

Generalization is implemented by
\[
 \mathsf{Gen}_i=
 \mathsf{Push};\mathsf{Offset}_i;\mathsf S;
        \mathsf A_{i,0,0,0};\mathsf S^{11}.
\]
It sends $(a,w,0,0)$ to $(U_k(i,a),w,0,0)$ for every $a,w,i$.
The Shift after Offset sees two visible zeros.  The saved old $a$ becomes
AX's cached input, and consequently its generalization head is exactly
$U_k(i,a)$.  The calculation remains valid for $a=1$.

\subsection{Compiler for every carried store}

Starting with $(0,[a,I_k(a,b),0;w],0,0)$, the word
\[
              \mathsf S;\mathsf S;\mathsf{MP}_b;\mathsf{Pop}
                                                               \tag{MPs}
\]
halts if the second Shift sees $a=1$; otherwise it returns
$(b,w,0,0)$.  The reserved zero below the two premises is necessary:
successful MP no longer pushes a dummy zero.  After successful MP the
stack is $[0;w]$ and the final Pop clears that entry while preserving
$b$, even when $b=1$.  If $a=b=w=0$, all zero aliases have the same
effect by total projection identities.

Induct on a finite numerical-calculus proof tree $e$ to construct a word
$W_e$ with the following exact specification: from every $(t,w,0,0)$ it
either halts or returns $(\operatorname{concl}(e),w,0,0)$.
For axiom leaves use (Sel) or the restricted selector.  For a
generalization node append $\mathsf{Gen}_i$.  For modus ponens with
major proof $e_1$ of $I_k(a,b)$ and minor proof $e_0$ of $a$, use
\[
 \mathsf{Discard};\mathsf{Push};W_{e_1};\mathsf{Push};
 W_{e_0};\mathsf{Push};\mathsf S;\mathsf S;
                     \mathsf{MP}_b;\mathsf{Pop}.          \tag{Compile}
\]
Discard and the first Push reserve zero.  The two inductive subwords
preserve their entire protected stack, and (MPs) supplies the inductive
step.  An earlier halt satisfies the specification rather than being
excluded by an unspoken consistency assumption.

If $e$ proves code one, then
\[
                         A_0;W_e;\mathsf S              \tag{HaltWord}
\]
halts from every store $(t,w,d,h)$.  The first AX clears offset and cache,
and if no previous halt occurs the last Shift sees $(1,0)$.
Appending Pop instead of Shift would be incorrect, since Pop is safe at
a visible one.

\subsection{Four-field enumeration and both implications}

Let $e(\epsilon)=0$ and $e(c::w)=1+b_c(e(w))$, where
\[
 \begin{aligned}
 b_{\mathsf S}(r)&=p(0,r),\\
 b_{\mathsf{MP}_x}(r)&=p(1,p(x,r)),\\
 b_{\mathsf A_{c,x,y,z}}(r)&=p(c+2,p(y,p(x,p(z,r)))).
 \end{aligned}                                           \tag{Encode}
\]
The AX fields occur in the order $y,x,z$ after the command tag.
The fourth field cannot be suppressed in the scheduler proof: a
three-field encoding makes the MP reader consume part of the continuation.
All inverse identities include zero fields and the zero Cantor alias.
The remainder after a command is a nested Cantor component, so it is at
most the code before the positive predecessor test and strictly below
the positive continuation code.  Hence every nonhalting candidate ends
after finitely many commands.  Each arithmetic macro has a terminating
all-natural contract, so this argument concerns actual finite source runs.

The fresh candidates increase successively from two.  For a code $r$ of
(HaltWord), prefixing $\mathsf A_{m,0,0,0}$ produces the unbounded family
of fresh codes
\[
                    p(m+2,p(0,p(0,p(0,r)))).
\]
The prefix may change the store; (HaltWord) works from every store.
Thus a numerical refutation supplies a reached halting candidate or the
machine halts earlier.  No reset between candidate words is assumed.

Conversely annotate $t$, $h$, and each live canonical stack entry by a
numerical-calculus derivation.  The initial zeros have fixed proofs.
AX creates only its ten theorem entries and saves a theorem top; this is
true for every fourth operand.  MP either applies modus ponens to two
annotated theorems or returns truth, and Shift moves annotations.  If the
sum guard holds, nonnegativity yields $(t,h)=(1,0)$ or $(0,1)$, so an
annotated contradiction is available.  Every finite canonical stack has
a primitive-recursive representation, including virtual zero entries.

The compiler, code packing, annotation transformations, and finite-run
arguments are primitive-recursive constructions whose correctness is
proved by ordinary induction in PA.  Together with (Cal), they establish
\[
 \mathrm{PA}\vdash
 H_{\mathrm{AX11Q6\ source}}
 \quad\longleftrightarrow\quad
 \exists p\,\operatorname{Prf}_{\mathrm{PA}}(p,\ulcorner0=1\urcorner).
\]
To infer the statement for the delivered literal TM, one additionally
needs its complete source/table binding, universal physical instruction
contracts with positive durations, and whole-tape blank initialization.
These must be replayed on the simultaneously modified final table; success
of three parent machines is not enough.  With those supplied checks and
their stated mathematical loop invariants, finite deterministic trace
composition transports the displayed equivalence to the literal machine.
Negating the existential proof predicate gives the desired PA-provable
nonhalting equivalence.  Neither this equivalence nor finite simulation
asserts that PA proves its own consistency.
